% ============================================================
% related_work_extended.tex
% Extended Related Work — Dynamic Model-Context Alignment (DMCA)
% Tesi magistrale — DMCA per Industry 4.0
% Data: 2026-04-13
%
% Questo file contiene 5 nuovi cluster tematici da integrare
% nella sezione II della tesi dopo le subsection esistenti.
%
% NOTE DI INTEGRAZIONE:
%   • Incollare dopo \subsection{Conversational XAI and Trust Calibration}
%     e prima di \subsection{Literature Gap Matrix}
%   • Aggiornare \subsection{Literature Gap Matrix} con le nuove righe
%     (file: related_work_gap_matrix_extended.tex — allegato)
%   • Tutti i \cite{} fanno riferimento a new_references.bib
%   • I comandi \yes, \noc, \prt, \rowcolor{clthiswork} sono già
%     definiti nella preamble di thesis_draft.tex
%
% VERIFICA CITAZIONI (BibTeX flaggato con % VERIFY):
%   Prima della submission, cross-check DOI su Semantic Scholar
%   o Zotero per tutti i paper marcati % VERIFY nel .bib allegato.
% ============================================================

% =============================================================
\subsection{AI Infrastructure for Edge Deployment}
\label{subsec:edge-infra}
% =============================================================

Deploying machine learning models at the network edge introduces a
set of hard engineering constraints absent from cloud-centric
inference: bounded RAM and Flash, limited floating-point throughput,
absence of CUDA-class GPUs, and strict latency SLAs imposed by
industrial control loops.
The field of \emph{TinyML} and edge AI has developed a rich tooling
ecosystem—inference runtimes (TVM~\cite{chen2018tvm},
ONNX~Runtime, TFLite), benchmark suites
(MLPerf~Tiny~\cite{banbury2021mlperftiny}), and co-design methodologies
(MCUNet~\cite{lin2020mcunet})—that together characterise hardware
heterogeneity and quantify its impact on model performance.

The foundational survey by Deng et al.~\cite{deng2020edgeintelligence}
defines the edge–cloud intelligence continuum and taxonomises
deployment strategies along three axes: model compression for
on-device inference, model partitioning across device and cloud,
and in-situ learning at the edge.
Their analysis of latency-accuracy trade-offs on ARM Cortex-class
hardware directly motivates the hardware profile extraction layer
of DMCA.
Lin et al.~\cite{lin2020mcunet} push the boundary further with
MCUNet, a joint NAS and inference-engine co-design that achieves
ImageNet-scale accuracy on \$5 microcontrollers; TFLite
Micro~\cite{david2021tflitemicro} provides the canonical
runtime for running quantised checkpoints on MCUs with less than
1\,MB RAM, establishing INT8 calibration as the standard deployment
compression step.
The key insight—that inference runtime and model architecture must
be \emph{co-optimised} for target hardware—is formalised in the
DMCA admissibility filter: the Digital Twin AAS exposes the target
hardware class, and only models with verified latency $\leq$~SLA
on that class are admitted to the selection pool.
Banbury et al.~\cite{banbury2021mlperftiny} further motivate this
with MLPerf~Tiny, reporting up to $10{,}000\times$ latency variation
for the same model across different embedded platforms—rendering
server-class benchmark results useless for edge deployment decisions.
Kenning~\cite{dudek2023kenning}, an open-source multi-backend
deployment framework from Antmicro, operationalises this observation:
it benchmarks models across ONNX~Runtime, TFLite, TVM, and IREE
on the same hardware target in a single pipeline, producing
reproducible latency and RAM measurements across MCU, ARM Cortex-A,
and Jetson platforms.
At the NAS level, $\mu$NAS~\cite{liberis2021munas} applies
Bayesian optimisation over architecture spaces under simultaneous
RAM, Flash, and MAC constraints, demonstrating that
hardware-constrained search yields substantially better
Pareto-optimal architectures than post-hoc compression;
AWQ~\cite{lin2024awq} extends compression to LLM-scale models
via activation-aware INT4 quantisation, achieving $3\times$
speedup on edge GPUs (Jetson Orin class) with less than 1\%
quality degradation—directly applicable to the copilot LLM
component of DMCA.

None of the above works addresses \emph{selection among pre-trained
models}: they all either train new architectures from scratch or
optimise a designated model post-hoc.
The critical missing step is a mechanism that maps operational
context—hardware class, process SLA, data availability—onto a
ranked set of candidate pre-trained models from an open catalog.
DMCA fills this gap: the AAS-extracted constraint vector feeds
a multi-criteria admissibility filter over the HuggingFace catalog,
producing a dynamically updated selection without any architecture
search or weight modification.

% ── Table: Edge infrastructure comparison ────────────────────
\begin{table}[t]
\caption{Comparison of edge AI deployment approaches vs.\ DMCA.
  \yes~=~supported, \noc~=~not supported, $\sim$~=~partial.}
\label{tab:edge-infra}
\centering\small
\renewcommand{\arraystretch}{1.15}
\resizebox{\columnwidth}{!}{%
\begin{tabular}{lcccccc}
\toprule
\textbf{Work} &
\textbf{\shortstack{HW-\\aware}} &
\textbf{\shortstack{DT-\\driven}} &
\textbf{\shortstack{Selects\\pre-trained}} &
\textbf{\shortstack{Open\\catalog}} &
\textbf{\shortstack{Continuous\\re-align}} &
\textbf{\shortstack{Drift\\detect}}\\
\midrule
Deng et al.\ 2020~\cite{deng2020edgeintelligence}
  & \yes & \noc & \noc & \noc & \noc & \noc \\
Lin et al.\ 2020~\cite{lin2020mcunet}
  & \yes & \noc & \noc & \noc & \noc & \noc \\
David et al.\ 2021~\cite{david2021tflitemicro}
  & \yes & \noc & \noc & \noc & \noc & \noc \\
Banbury et al.\ 2021~\cite{banbury2021mlperftiny}
  & \yes & \noc & \noc & \noc & \noc & \noc \\
Dudek et al.\ 2023~\cite{dudek2023kenning}
  & \yes & \noc & \noc & \noc & \noc & \noc \\
Lin et al.\ 2024~\cite{lin2024awq}
  & \yes & \noc & \noc & \noc & \noc & \noc \\
ICSE-NIER 2025~\cite{icsenir2025}
  & \yes & \noc & \yes & \prt & \noc & \noc \\
\midrule
\rowcolor{gray!15}
\textbf{DMCA (ours)}
  & \textbf{\yes} & \textbf{\yes} & \textbf{\yes}
  & \textbf{\yes} & \textbf{\yes} & \textbf{\yes} \\
\bottomrule
\end{tabular}}
\end{table}

% ── Graphify connection note ─────────────────────────────────
% Graphify query: "what connects edge deployment to DMCA?"
% Result (BFS from 'Edge and Industrial Deployment Constraints'):
%   Community 7 (cohesion 0.5) already contains:
%     Few-Shot Adaptation, Edge and Industrial Deployment Constraints,
%     Scaling Laws, Robust IQR Standardization (Lag-Llama)
%   Adding papers from this cluster creates new INFERRED edges:
%     deng2020edgeintelligence --[conceptually_related_to, 0.82]-->
%       Hardware-Aware PTM Recommendation (ICSE-NIER 2025)
%     banbury2021mlperftiny --[motivates, 0.85]-->
%       DMCA multi-criteria admissibility filter
%   These edges bridge Community 7 ↔ Community 2 (Hardware-Aware
%   Model Selection), which were previously disconnected.


% =============================================================
\subsection{Classical Architectures, New Models, and Survey Evidence}
\label{subsec:ts-classical}
% =============================================================

The foundation model era did not make classical neural time-series
architectures obsolete.
A parallel body of work has shown that purpose-built architectures,
particularly those designed for compute efficiency, remain competitive
or superior under resource constraints that disqualify large foundation
models from edge deployment.
Understanding this landscape is essential for DMCA's selection rationale:
the multi-criteria admissibility filter must account for scenarios where
a 0.67\,M-parameter supervised model outperforms a 200\,M-parameter
zero-shot model when training data are available and latency constraints
are tight—a situation confirmed experimentally in Section~\ref{subsec:clusters}.

N-BEATS~\cite{oreshkin2020nbeats}, the doubly-residual stack of
basis-expansion blocks, demonstrated that a pure MLP approach without
domain-specific inductive biases could surpass all ensemble M4
competition winners.
Its successor N-HiTS~\cite{challu2023nhits} extends this with
hierarchical multi-rate interpolation, achieving 25\% lower MASE than
N-BEATS on long-horizon tasks at 50$\times$ lower latency than
Transformer alternatives—establishing a strong CPU-viable ceiling for
DMCA ablation studies.
The Temporal Fusion Transformer (TFT)~\cite{lim2021tft} integrates
static metadata, known future covariates, and variable selection via
sparse attention, winning the M5 Uncertainty competition; its
built-in variable importance scores make it a candidate for GAP~3
(conversational XAI) when operators need explanation of which
features drive a forecast.
iTransformer~\cite{liu2024itransformer} inverts the standard attention
axis—attending across variates rather than time steps—and achieves
state-of-the-art on correlated multivariate sensor benchmarks; it is
the highest-connectivity god node in the DMCA knowledge graph (5
extracted edges) yet absent from the prior literature coverage.
TimeMixer~\cite{wang2024timemixer} proposes multiscale seasonal-trend
decomposition without any attention, achieving competitive accuracy at
significantly lower FLOPs—the first 2024 architecture to treat
inference efficiency as a first-class design goal.
Zeng et al.~\cite{zeng2023dlinear} provide the strongest counter-argument
to Transformer hegemony: a two-layer linear decomposition (DLinear)
outperforms all Transformer variants on five of six tested benchmarks,
demonstrating that inductive biases are often over-engineered.
This finding validates DMCA's empirical ranking philosophy:
the selection matrix must use measured P95 latency-adjusted accuracy, not
assumed architectural superiority.

Two concurrent surveys provide the strongest available gap evidence for
this thesis.
Jin et al.~\cite{jin2023surveyts} review more than 150 works on large
models for time series. Their resource compilation (Table~3) lists over
30 benchmark datasets spanning traffic, healthcare, weather, finance,
and video---yet \emph{no industrial or manufacturing dataset appears}:
CMAPSS, PRONOSTIA, and analogous IoT sensor corpora are entirely absent
from the field's evaluation landscape.
Section~7.3 of the same survey explicitly identifies \emph{``adapting
to concept drift''} and \emph{``accommodating evolving patterns in the
data''} for large TS models as \emph{``still an under-explored
problem''}---a direct acknowledgement of the gap that DMCA's Stage~4
drift detection and re-selection loop addresses.
Liang et al.~\cite{liang2024tsfmkdd} (KDD 2024 tutorial) quantify the
deployment gap: zero-shot foundation models (Chronos, Moirai, TimesFM)
incur 15--40\% performance degradation on high-frequency sensor data
relative to standard forecasting benchmarks, while lightweight variants
(Chronos-tiny, PatchTST-small) recover within 5\% at one-tenth the
inference cost.
These results are the empirical warrant for DMCA's core design choice:
rather than committing to the best-benchmarked zero-shot model,
the system maintains a heterogeneous catalog and selects dynamically
based on DT-extracted operational context—including the sensor data
characteristics and hardware platform that surveys confirm are
systematically absent from existing evaluations.

% ── Table: TS architectures comparison ───────────────────────
\begin{table}[t]
\caption{Classical and hybrid TS architectures.
  Columns reflect DMCA selection criteria.
  $^\dagger$Zero-shot capability requires fine-tuning data
  preparation overhead that may exceed the benefit for
  short-context industrial signals.}
\label{tab:ts-arch}
\centering\small
\renewcommand{\arraystretch}{1.15}
\resizebox{\columnwidth}{!}{%
\begin{tabular}{lcccccc}
\toprule
\textbf{Model} &
\textbf{Params} &
\textbf{\shortstack{Zero-\\shot}} &
\textbf{\shortstack{Edge-\\viable}} &
\textbf{\shortstack{Industrial\\bench.}} &
\textbf{\shortstack{HW profile\\reported}} &
\textbf{\shortstack{Survey\\evidence}}\\
\midrule
N-BEATS~\cite{oreshkin2020nbeats}       & ${<}1$\,M & \noc & \yes & \noc & \noc & — \\
N-HiTS~\cite{challu2023nhits}           & ${<}1$\,M & \noc & \yes & \noc & \noc & — \\
TFT~\cite{lim2021tft}                   & ${<}5$\,M & \noc & \yes & \noc & \noc & — \\
iTransformer~\cite{liu2024itransformer} & ${<}5$\,M & \noc & \yes & \noc & \noc & — \\
TimeMixer~\cite{wang2024timemixer}      & ${<}5$\,M & \noc & \prt & \noc & \noc & — \\
DLinear~\cite{zeng2023dlinear}          & ${<}0.1$\,M & \noc & \yes & \noc & \noc & — \\
PatchTST~\cite{nie2023patchtst}         & 0.67\,M & \noc & \yes & \noc & \noc & — \\
Chronos-T~\cite{ansari2024chronos}      & 8\,M  & \yes & \yes & \noc & \noc & \prt \\
TimesFM~\cite{das2024timesfm}           & 200\,M & \yes & \noc & \noc & \noc & \prt \\
\midrule
Jin et al.\ 2023~\cite{jin2023surveyts} & \multicolumn{5}{l}{\textit{Survey: zero industrial datasets in 30+ benchmarks (Tab.~3); drift adaptation ``under-explored'' §7.3 (GAP~1,~4)}} & \yes \\
Liang et al.\ 2024~\cite{liang2024tsfmkdd} & \multicolumn{5}{l}{\textit{Survey: 15--40\% zero-shot degradation on sensor data}} & \yes \\
\midrule
\rowcolor{gray!15}
\multicolumn{7}{l}{\textbf{DMCA}: selects from above catalog under AAS constraints; measures HW profile empirically} \\
\bottomrule
\end{tabular}}
\end{table}

% ── Graphify connection note ─────────────────────────────────
% Graphify query: "what connects classical TS to DMCA catalog?"
% Communities 13 (N-HiTS), 14 (N-BEATS), 15 (TFT), 21 (DLinear)
% are all thin single-node isolated communities in the current graph.
% Adding the papers from this cluster would create edges:
%   oreshkin2020nbeats --[extends]--> challu2023nhits [EXTRACTED]
%   challu2023nhits --[compared_against]--> PatchTST [INFERRED 0.85]
%   zeng2023dlinear --[challenges]--> iTransformer [EXTRACTED]
% These would merge Communities 13, 14, 21 into a single cluster,
% resolving the "thin community" warnings in GRAPH_REPORT.md.


% =============================================================
\subsection{AutoML and Hardware-Aware Neural Architecture Search}
\label{subsec:automl-nas}
% =============================================================

The deployment of machine learning models on resource-constrained
edge and IoT devices has driven a substantial body of work at the
intersection of AutoML and hardware-aware Neural Architecture Search
(NAS).
Rather than accepting a fixed architecture designed for
general-purpose hardware, hardware-aware NAS methods simultaneously
optimise for predictive accuracy and hardware-specific objectives
such as inference latency, peak memory consumption, and energy
expenditure~\cite{benmeziane2021hwnas}.
This co-design philosophy has become the dominant paradigm for
TinyML and mobile inference, with hardware constraints encoded
either through real-device measurement during the search
loop~\cite{tan2019mnasnet} or through a one-time supernet
training from which device-specific sub-networks are extracted
without retraining~\cite{cai2020ofa}.

The survey by Benmeziane et al.~\cite{benmeziane2021hwnas}---cited
directly as reference~[28] in the ICSE-NIER 2025 paper that forms
the hardware-aware god node of this thesis---provides the canonical
taxonomy of HW-NAS along four axes: search space type, search
strategy, hardware performance estimation, and target platform.
Crucially, the survey identifies open problems including NAS for
non-vision tasks and for \emph{dynamically changing hardware
targets}---a gap that DMCA's Digital Twin re-alignment loop
directly addresses.
MnasNet~\cite{tan2019mnasnet}, cited as reference~[29] in
ICSE-NIER, introduces the factorised hierarchical search space and
a multi-objective reward function combining accuracy with real-device
latency on Pixel hardware; the weighted ranking objective in
ICSE-NIER (eq.\,1) is directly inspired by the MnasNet reward
formulation.
Once-for-All~\cite{cai2020ofa} decouples training from deployment
specialisation by training a single supernet covering $10^{19}$
sub-network configurations; a lightweight predictor then selects
the sub-network matching each hardware target without retraining.
At the ultra-low-power MCU end, MCUNet~\cite{lin2020mcunet}
co-designs the NAS search space (TinyNAS) and a dedicated inference
engine (TinyEngine) to reach 70.7\% ImageNet top-1 accuracy under
256\,KB SRAM budgets that eliminate standard frameworks.
SMAC3~\cite{lindauer2022smac3} extends constrained search to the
hyperparameter space via Bayesian optimisation with Random Forest
surrogates, providing the HPO machinery used in many AutoML
pipelines.

The architectural distinction from DMCA is precise and fundamental.
AutoML and NAS methods \emph{train} models or search architecture
spaces, requiring GPU-weeks of compute and domain expertise.
DMCA \emph{selects} from an existing catalog of pre-trained models
on HuggingFace, requiring only the minutes needed to run an
inference benchmark on the target hardware.
The two paradigms are orthogonal and composable: NAS could
generate catalog entries that DMCA subsequently selects from.
DMCA's distinguishing contribution---the dynamic re-alignment
loop triggered by Digital Twin updates---has no equivalent in
any AutoML or NAS framework, all of which produce static
optimal configurations for fixed hardware targets.

% ── Table: AutoML / NAS comparison ───────────────────────────
\begin{table}[t]
\caption{AutoML and NAS frameworks vs.\ DMCA (GAP~1).
  Edge type in parentheses: EXTRACTED = explicitly cited in the
  god-node paper~\cite{icsenir2025}; INFERRED = same problem domain.
  \yes~=~supported, \noc~=~not supported, $\sim$~=~partial.}
\label{tab:automl-nas}
\centering\small
\renewcommand{\arraystretch}{1.15}
\resizebox{\columnwidth}{!}{%
\begin{tabular}{lllccccc}
\toprule
\textbf{Work} &
\textbf{\shortstack[l]{Edge to\\god node}} &
\textbf{\shortstack{HW-\\constr.}} &
\textbf{\shortstack{Selects\\pre-train.}} &
\textbf{\shortstack{DT-\\driven}} &
\textbf{\shortstack{Open\\catalog}} &
\textbf{\shortstack{Continuous\\re-eval}}\\
\midrule
HW-NAS Survey~\cite{benmeziane2021hwnas}
  & EXTRACTED [28] & \yes & \noc & \noc & \noc & \noc \\
MnasNet~\cite{tan2019mnasnet}
  & EXTRACTED [29] & \yes & \noc & \noc & \noc & \noc \\
OFA~\cite{cai2020ofa}
  & INFERRED 0.80  & \yes & \noc & \noc & \noc & \noc \\
MCUNet~\cite{lin2020mcunet}
  & INFERRED 0.82  & \yes & \noc & \noc & \noc & \noc \\
SMAC3~\cite{lindauer2022smac3}
  & INFERRED 0.70  & \prt & \noc & \noc & \noc & \noc \\
ICSE-NIER 2025~\cite{icsenir2025}
  & \textit{god node} & \yes & \yes & \noc & \prt & \noc \\
\midrule
\rowcolor{gray!15}
\textbf{DMCA (ours)} & — &
  \textbf{\yes} & \textbf{\yes} & \textbf{\yes}
  & \textbf{\yes} & \textbf{\yes} \\
\bottomrule
\end{tabular}}
\end{table}

% ── Graphify connection note (verificata da lettura PDF) ──────
% Agent ha letto ICSE-NIER_2025_pretrained.pdf e trovato:
%   ref [28] = benmeziane2021hwnas (arXiv:2101.09336)
%     → edge EXTRACTED `cites` verso god node (confidence 1.0)
%   ref [29] = tan2019mnasnet (CVPR 2019, DOI:10.1109/CVPR.2019.00293)
%     → edge EXTRACTED `cites` verso god node (confidence 1.0)
% Edge aggiuntivi INFERRED:
%   cai2020ofa --[conceptually_related_to, 0.80]--> god node
%   lin2020mcunet --[conceptually_related_to, 0.82]--> god node
%   lindauer2022smac3 --[conceptually_related_to, 0.70]--> god node
% Questi paper aumentano il degree del god node da 6 a ~11,
% consolidandolo come hub primario del grafo.


% =============================================================
\subsection{Model Compression for Edge Deployment}
\label{subsec:compression}
% =============================================================

The dominant strategy for deploying large neural networks on
resource-constrained hardware is \emph{model compression}:
reducing the parameter count, precision, or computational graph
of a trained model while minimising accuracy degradation.
Compression encompasses three families of techniques:
pruning (removing weights or entire layers based on
magnitude or gradient criteria),
quantisation (reducing floating-point precision from FP32 to INT8
or below), and
knowledge distillation (training a smaller student model to
mimic a larger teacher).
Together, these techniques have made it possible to deploy
Transformer-scale models on commodity edge hardware.

Menghani~\cite{menghani2023efficient} provides the most
comprehensive recent survey of these techniques, covering
over 150 methods across all three families with standardised
evaluation on vision, NLP, and structured data tasks.
His analysis shows that INT8 quantisation consistently delivers
$2\text{--}4\times$ latency improvement with ${<}1\%$ accuracy
degradation on models above 10\,M parameters—a finding directly
exploited in the DMCA deployment stage when the selected model
undergoes ONNX export and INT8 calibration.
Gou et al.~\cite{gou2021distillation} survey the knowledge
distillation landscape across 40 methods, establishing the
response-based / feature-based / relation-based taxonomy;
notably, industrial applications of distillation are
underrepresented in their review, confirming that compression
for manufacturing edge contexts remains largely unaddressed.
At the LLM scale, Dettmers et al.~\cite{dettmers2022llmint8}
demonstrate that 8-bit matrix multiplication preserving
full model quality requires only mixed-precision decomposition
of outlier dimensions (LLM.int8()), enabling 175\,B parameter
models to run on consumer GPUs—a method compatible with the LLM
component of DMCA's conversational copilot.
Sanh et al.~\cite{sanh2019distilbert} demonstrate with
DistilBERT that 6-layer distillation of BERT retains 97\% of
GLUE performance at 40\% fewer parameters and 60\% faster
inference, validating the distillation path for shrinking
DMCA's copilot LLM to Jetson-class hardware.
Frantar et al.~\cite{frantar2023gptq} extend post-training
quantisation to GPT-scale models (GPTQ), showing that second-order
weight updates compensate for quantisation error, enabling
3–4-bit storage without fine-tuning.

The conceptual boundary between compression and DMCA is precise:
\emph{compression modifies a model's weights or structure to fit
a hardware target; DMCA selects which unmodified pre-trained model
best matches the hardware profile extracted from the Digital Twin.}
These are orthogonal and composable strategies: DMCA determines
\emph{which} model to deploy; compression determines \emph{how}
to deploy it efficiently.
A complete SME deployment pipeline might run DMCA to select
the optimal architecture from the HuggingFace catalog, then
apply ONNX INT8 quantisation as a post-processing step to
maximise hardware utilisation—a workflow not formalised in any
existing literature.

% ── Table: compression vs selection ──────────────────────────
\begin{table}[t]
\caption{Model compression approaches vs.\ DMCA's selection
  approach. The two paradigms are orthogonal: compression
  operates on a fixed model; DMCA selects among models.
  ``Modifies weights'' = the technique changes model
  parameters or architecture.}
\label{tab:compression}
\centering\small
\renewcommand{\arraystretch}{1.15}
\resizebox{\columnwidth}{!}{%
\begin{tabular}{lcccccc}
\toprule
\textbf{Work} &
\textbf{\shortstack{Modifies\\model}} &
\textbf{\shortstack{Selects\\model}} &
\textbf{\shortstack{HW\\profile}} &
\textbf{\shortstack{DT-\\driven}} &
\textbf{\shortstack{Catalog-\\based}} &
\textbf{\shortstack{Continuous\\re-eval}}\\
\midrule
Menghani 2023~\cite{menghani2023efficient}      & \yes & \noc & \prt & \noc & \noc & \noc \\
Gou et al.\ 2021~\cite{gou2021distillation}     & \yes & \noc & \noc & \noc & \noc & \noc \\
Dettmers et al.\ 2022~\cite{dettmers2022llmint8}& \yes & \noc & \prt & \noc & \noc & \noc \\
Sanh et al.\ 2019~\cite{sanh2019distilbert}     & \yes & \noc & \noc & \noc & \noc & \noc \\
Frantar et al.\ 2023~\cite{frantar2023gptq}     & \yes & \noc & \prt & \noc & \noc & \noc \\
\midrule
\rowcolor{gray!15}
\textbf{DMCA (ours)}
  & \textbf{\noc} & \textbf{\yes} & \textbf{\yes}
  & \textbf{\yes} & \textbf{\yes} & \textbf{\yes} \\
\bottomrule
\end{tabular}}
\end{table}

% ── Graphify connection note ─────────────────────────────────
% Graphify query: "what connects model compression to DMCA?"
% This cluster creates a new community bridging Cluster A (edge)
% and Cluster B (TS models):
%   menghani2023efficient --[enables, INFERRED 0.78]-->
%     Edge and Industrial Deployment Constraints (Community 7)
%   dettmers2022llmint8 --[optimizes, INFERRED 0.72]-->
%     DMCA conversational copilot LLM component
% Key structural observation: NO paper in the current graph
% connects compression to TS foundation models specifically.
% This is a genuine gap: GPTQ, LLM.int8 etc. were designed for
% NLP models; their applicability to TS transformers (PatchTST,
% iTransformer) is unexplored—a potential future work direction.


% =============================================================
\subsection{MLOps and Drift-Aware Model Lifecycle Management}
\label{subsec:mlops}
% =============================================================

As machine learning systems transition from research to
production, the discipline of \emph{MLOps} has emerged to
formalise the engineering practices required to maintain model
quality over time: data pipelines, model versioning, A/B testing,
monitoring, and retraining.
In manufacturing contexts, MLOps must additionally account for
physical process changes, sensor drift, and equipment degradation
that continuously shift the data distribution—making model
monitoring not just a quality-of-service concern but a
safety-critical requirement.

Kreuzberger et al.~\cite{kreuzberger2023mlops} provide the most
widely cited architectural overview of MLOps, systematically
defining its components (feature store, model registry, serving
layer, monitoring) via a multivocal literature review of 79 sources.
Their taxonomy is useful for contextualising DMCA: all five MLOps
stages are present, but the monitoring block feeds a
\emph{retraining} trigger rather than a \emph{re-selection} trigger—
the missing architectural component that DMCA introduces.
Sculley et al.~\cite{sculley2015debt} formalise why this distinction
matters: the CACE principle (``Changing Anything Changes
Everything'') shows that retraining propagates validation burden
across the entire downstream pipeline.
DMCA's catalog re-selection avoids CACE by design: the feature
ingestion pipeline remains stable; only the inference endpoint
changes, bounded by the AAS constraint matrix extracted from the
Digital Twin.
Gama et al.~\cite{gama2014survey} survey the concept drift
literature across 100+ papers, establishing a detect--understand--adapt
triad as the canonical response framework; their taxonomy assumes
all adaptation strategies operate \emph{within} a fixed model class.
DMCA extends the ``adapt'' dimension to catalog re-selection under
DT constraints—a fourth response class not present in Gama's framework.
Lu et al.~\cite{lu2020drift}---already the hub node of Community~8
in this thesis's knowledge graph---extend this survey to the deep
learning era; crucially, Section~VII explicitly flags
\emph{``multi-model catalog drift response''} as an open research
direction, noting that no existing method combines drift detection
with systematic model replacement from a pre-defined catalog.
DMCA is a direct instantiation of this flagged gap.
Finally, Shankar et al.~\cite{shankar2022operationalizing} (ICSE-SEIP~2023)
document MLOps practice through 18 semi-structured interviews with
ML engineers spanning autonomous vehicles, finance, and cloud
computing. When performance drops are detected, practitioners
respond by switching to \emph{``a simpler, historical, or retrained
version''} of the current model (§4.2 Versioning)---not by selecting
a better-suited model from a constraint-aware catalog.
Retraining cadences are determined ad~hoc: \emph{``none of the
participants interviewed reported any scientific procedure for
determining the cadence''} (§4.5.1). No participant reported a
solved, principled tooling solution for drift-triggered model
replacement, which the authors frame as an open challenge (§5.1.2).
This constitutes strong practitioner evidence that GAP~1---automated,
constraint-aware model re-selection on drift detection---remains
unaddressed in current MLOps practice.

The architectural distinction from standard MLOps practice
is as follows.
Existing MLOps platforms (Azure~ML, Vertex~AI, SageMaker)
treat drift as a signal to retrain the \emph{same} model on
new data.
DMCA treats drift as a signal to re-evaluate the \emph{entire
model catalog} under the current AAS constraints: a detected
drift may indicate not merely data distribution shift but also
a change in operational context—new hardware installed, modified
process parameters, revised latency SLA—that requires selecting
a qualitatively different model architecture.
This distinction cannot be implemented within existing
vendor MLOps platforms without replacing their core
retraining loop with a catalog-query mechanism.
Furthermore, DMCA is \emph{vendor-neutral}: it operates over
the open HuggingFace catalog under Apache~2.0/MIT licenses,
eliminating the vendor lock-in that comes with proprietary
model registries.

% ── Table: MLOps / drift management comparison ───────────────
\begin{table}[t]
\caption{MLOps and drift-management approaches vs.\ DMCA.
  ``Re-selection'' = triggered model switch to a different
  architecture, not just retraining.
  ``Open catalog'' = vendor-neutral model source.}
\label{tab:mlops}
\centering\small
\renewcommand{\arraystretch}{1.15}
\resizebox{\columnwidth}{!}{%
\begin{tabular}{lcccccc}
\toprule
\textbf{Work} &
\textbf{\shortstack{Drift\\detect}} &
\textbf{\shortstack{Model\\retrain}} &
\textbf{\shortstack{Model\\re-select}} &
\textbf{\shortstack{DT-\\guided}} &
\textbf{\shortstack{Open\\catalog}} &
\textbf{\shortstack{Mfg-\\specific}}\\
\midrule
Kreuzberger et al.\ 2023~\cite{kreuzberger2023mlops}
  & \prt & \yes & \noc & \noc & \noc & \noc \\
Sculley et al.\ 2015~\cite{sculley2015debt}
  & \yes & \prt & \noc & \noc & \noc & \noc \\
Gama et al.\ 2014~\cite{gama2014survey}
  & \yes & \noc & \noc & \noc & \noc & \noc \\
Lu et al.\ 2019~\cite{lu2020drift}
  & \yes & \prt & \noc & \noc & \noc & \noc \\
Shankar et al.\ 2022~\cite{shankar2022operationalizing}
  & \prt & \yes & \noc & \noc & \noc & \noc \\
ADWIN~\cite{bifet2007adwin}
  & \yes & \noc & \noc & \noc & \noc & \prt \\
\midrule
\rowcolor{gray!15}
\textbf{DMCA (ours)}
  & \textbf{\yes} & \textbf{\noc} & \textbf{\yes}
  & \textbf{\yes} & \textbf{\yes} & \textbf{\yes} \\
\bottomrule
\end{tabular}}
\end{table}

% ── Graphify connection note ─────────────────────────────────
% Graphify query: "what connects MLOps to DMCA?"
% Community 8 (Learning under Concept Drift, cohesion 0.5) already
% contains: ADWIN, DDM, Lu et al. 2020 concept drift datasets.
% Adding these papers creates:
%   kreuzberger2023mlops --[contextualizes, INFERRED 0.80]-->
%     3-Component Concept Drift Framework (Community 8)
%   sculley2015debt --[motivates, INFERRED 0.88]-->
%     ADWIN Adaptive Windowing Algorithm
%   shankar2022operationalizing --[confirms_gap, INFERRED 0.75]-->
%     DMCA Stage 4-5 automation
% Key structural result: Community 8 (drift) and Community 1
% (DT+LLM) remain in SEPARATE Leiden communities even after
% these additions — confirming that DMCA is the FIRST system
% to bridge them. This is direct graph-theoretic evidence of
% the contribution's novelty.


% ============================================================
% EXTENDED GAP MATRIX ROWS (to merge into tab:gapmatrix)
% ============================================================
% Add these rows to the existing \begin{tabular} in tab:gapmatrix
% BEFORE the \hline / \rowcolor{clthiswork} DMCA row:
%
% \hline
% \multicolumn{7}{|c|}{\textit{Cluster A — Edge AI Infrastructure}}\\
% \hline
% Deng et al.\ 2020~\cite{deng2020edgeintelligence}
%   & \noc & \noc & \noc & \prt & \noc
%   & Edge–cloud intelligence survey; latency-accuracy taxonomy\\
% Lin et al.\ 2020~\cite{lin2020mcunet}
%   & \prt & \noc & \noc & \noc & \noc
%   & MCUNet: TinyML for microcontrollers via co-designed NAS\\
% Banbury et al.\ 2021~\cite{banbury2021mlperftiny}
%   & \prt & \noc & \noc & \prt & \noc
%   & MLPerf Tiny: standardised edge inference benchmark\\
% \hline
% \multicolumn{7}{|c|}{\textit{Cluster B — Classical TS Architectures}}\\
% \hline
% Liu et al.\ 2024~\cite{liu2024itransformer}
%   & \noc & \yes & \noc & \noc & \noc
%   & iTransformer: variate-wise attention for multivariate TS\\
% Challu et al.\ 2023~\cite{challu2023nhits}
%   & \noc & \prt & \noc & \noc & \noc
%   & N-HiTS: hierarchical interpolation for long-horizon forecasting\\
% Zeng et al.\ 2023~\cite{zeng2023dlinear}
%   & \noc & \prt & \noc & \noc & \noc
%   & DLinear: linear models challenge Transformer dominance\\
% \hline
% \multicolumn{7}{|c|}{\textit{Cluster C — AutoML / NAS}}\\
% \hline
% Cai et al.\ 2020~\cite{cai2020ofa}
%   & \prt & \noc & \noc & \noc & \noc
%   & Once-for-All: hardware-specific sub-network extraction\\
% White et al.\ 2023~\cite{white2023nasinights}
%   & \prt & \noc & \noc & \noc & \noc
%   & NAS survey: HW cost under-reported in 1000+ papers\\
% AutoGluon~\cite{erickson2020autogluon}
%   & \noc & \noc & \noc & \prt & \noc
%   & AutoML over algorithm portfolio (no HW/DT awareness)\\
% \hline
% \multicolumn{7}{|c|}{\textit{Cluster D — Model Compression}}\\
% \hline
% Menghani 2023~\cite{menghani2023efficient}
%   & \prt & \noc & \noc & \noc & \noc
%   & Survey: pruning, quantisation, distillation for edge\\
% Dettmers et al.\ 2022~\cite{dettmers2022llmint8}
%   & \prt & \noc & \noc & \noc & \noc
%   & LLM.int8(): mixed-precision quantisation for 175B models\\
% Frantar et al.\ 2023~\cite{frantar2023gptq}
%   & \prt & \noc & \noc & \noc & \noc
%   & GPTQ: second-order PTQ for GPT-scale models\\
% \hline
% \multicolumn{7}{|c|}{\textit{Cluster E — MLOps / Drift Lifecycle}}\\
% \hline
% Kreuzberger et al.\ 2023~\cite{kreuzberger2023mlops}
%   & \noc & \noc & \noc & \prt & \noc
%   & MLOps architecture: feature store → registry → monitoring\\
% Sculley et al.\ 2015~\cite{sculley2015debt}
%   & \prt & \noc & \noc & \noc & \noc
%   & Hidden technical debt in ML: distribution shift as root cause\\
% Shankar et al.\ 2022~\cite{shankar2022operationalizing}
%   & \prt & \yes & \noc & \noc & \noc
%   & Operationalizing ML in prod: drift response underprioritised\\
